% Part: computability
% Chapter: computability-theory
% Section: computable-sets

\documentclass[../../../include/open-logic-section]{subfiles}

\begin{document}

\olfileid[fr]{cmp}{thy}{cps}
\olsection{Ensembles calculables}

La notion de calculabilité s'étend des
fonctions aux ensembles :

\begin{defn}
  Soit $S$ un ensemble d'entiers naturels.
  L'ensemble $S$ est \emph{calculable} si et
  seulement si sa fonction caractéristique
  $\Char{S}$ l'est. Autrement dit, $S$ est
  calculable si et seulement si la fonction
\[
\Char{S}(x) =
\begin{cases}
1 & \text{si $x\in S$}\\
0 & \text{sinon}
\end{cases}
\]
  est calculable. De même, une relation
  $R(x_0,\dots,x_{k-1})$ est calculable si
  et seulement si sa fonction caractéristique
  l'est.

Les ensembles et les relations calculables
sont également dits \emph{décidables}.
\end{defn}

\begin{explain}
Nous disposons maintenant de plusieurs notions
de calculabilité : pour les fonctions partielles,
pour les fonctions totales et pour les ensembles.
Il faut les distinguer.\iftag{TMs}{ Une machine
de Turing qui calcule une fonction partielle
renvoie la valeur de cette fonction pour les
entrées où elle est définie. Une machine qui
décide un ensemble s'arrête pour toute entrée
et renvoie $1$ ou~$0$ selon que celle-ci
appartient ou non à l'ensemble.}{}
\end{explain}

\end{document}
